% GENERATED FILE. Edit canonical lessons, then run npm run generate:books.
% canonical-source-hash: fnv1a64:a876eea86f66e02c
% canonical-lessons: TE-C499-asuyapadu, TE-C499-dvesincu, TE-C499-asahyincuko, TE-C499-paga, TE-C499-cinnacupu

\chapter{Envy and Dislike}
\label{ch:te-envy-and-dislike}
\hlchaptermodality{499}

\begin{chapteropening}
\textbf{By the end of this chapter:} \emph{I can say that someone envied, hated or looked down on another.}

Complete the last lesson of chapter 499: I can say that someone envied, hated or looked down on another.
\end{chapteropening}

\section[\texorpdfstring{\te{అసూయపడు} (asūyapaḍu)}{అసూయపడు (asūyapaḍu)}]{\te{అసూయపడు} (asūyapaḍu) — to be jealous; to envy}

\label{lesson:TE-C499-asuyapadu}



\begin{quote}
Before the new one: say the Telugu for to hold back politely, then the Telugu for to hang one's head.
\end{quote}

\subsection*{You'll want to know: \te{అసూయపడు}}
\textbf{\te{అసూయపడు}} — \emph{asūyapaḍu} — \textquotedblleft{}to be jealous; to envy\textquotedblright{}.

\begin{quote}
\te{అతని} \te{కొత్త} \te{కారు} \te{చూసి} \te{స్నేహితులు} \te{అసూయపడ్డారు}.
\end{quote}

\begin{quote}
\emph{atani kotta kāru cūsi snēhitulu asūyapaḍḍāru.}
\end{quote}

\begin{quote}
Seeing his new car, his friends were jealous.
\end{quote}

Past (I …): \te{అసూయపడ్డాను}, \emph{asūyapaḍḍānu}.

Who you envy takes \te{మీద}: \te{అతని} \te{మీద} \te{అసూయపడు}.

\begin{cousinweb}
\textbf{\te{అసూయపడు}} (\emph{asūyapaḍu}) is built from \te{అసూయ}, \textquotedblleft{}envy\textquotedblright{}, from Sanskrit \emph{asūyā}, resentment at another person's good fortune (Monier-Williams) and the native verb \emph{paḍu}, \textquotedblleft{}to feel, to suffer\textquotedblright{} (DEDR 3853).
\end{cousinweb}

\begin{grammarlens}[title={What you've built}]
One more word for describing what you lived and felt.
\end{grammarlens}

\subsection*{Your turn}
\begin{itemize}
\raggedright
  \item \emph{Say it:} \emph{asūyapaḍu}
  \item \emph{Say it:} \emph{asūyapaḍu}, once more
  \item \emph{Say it:} \emph{atani kotta kāru cūsi snēhitulu asūyapaḍḍāru.} — Seeing his new car, his friends were jealous.
  \item \emph{Recall:} say the Telugu for to hold back politely, then the Telugu for to hang one's head, then say \emph{asūyapaḍu} again
\end{itemize}

\begin{tcolorbox}[breakable,colback=teal!4,colframe=teal!35!black,title={Before you move on}]
What does \te{అసూయపడు} mean? (\textquotedblleft{}To be jealous; to envy\textquotedblright{}.) Say it once more.
\end{tcolorbox}
\section[\texorpdfstring{\te{ద్వేషించు} (dvēṣiñcu)}{ద్వేషించు (dvēṣiñcu)}]{\te{ద్వేషించు} (dvēṣiñcu) — to hate}

\label{lesson:TE-C499-dvesincu}



\begin{quote}
Before the new one: say the Telugu for to hang one's head, then the Telugu for to be jealous.
\end{quote}

\subsection*{You'll want to know: \te{ద్వేషించు}}
\textbf{\te{ద్వేషించు}} — \emph{dvēṣiñcu} — \textquotedblleft{}to hate\textquotedblright{}.

\begin{quote}
\te{అబద్ధాలను} \te{నేను} \te{ద్వేషిస్తాను}.
\end{quote}

\begin{quote}
\emph{abaddhālanu nēnu dvēṣistānu.}
\end{quote}

\begin{quote}
I hate lies.
\end{quote}

Past (I …): \te{ద్వేషించాను}, \emph{dvēṣiñcānu}.

The noun \te{ద్వేషం} is hatred; its opposite is \te{ప్రేమ}.

\begin{cousinweb}
\textbf{\te{ద్వేషించు}} (\emph{dvēṣiñcu}) is built on Sanskrit \emph{dvēṣa}, \textquotedblleft{}hatred\textquotedblright{} (CDIAL 6696), from the root \emph{dviṣ}, with the Telugu verb ending -\te{ఇంచు}.
\end{cousinweb}

\begin{grammarlens}[title={What you've built}]
One more word for describing what you lived and felt.
\end{grammarlens}

\subsection*{Your turn}
\begin{itemize}
\raggedright
  \item \emph{Say it:} \emph{dvēṣiñcu}
  \item \emph{Say it:} \emph{dvēṣiñcu}, once more
  \item \emph{Say it:} \emph{abaddhālanu nēnu dvēṣistānu.} — I hate lies.
  \item \emph{Recall:} say the Telugu for to hang one's head, then the Telugu for to be jealous, then say \emph{dvēṣiñcu} again
\end{itemize}

\begin{tcolorbox}[breakable,colback=teal!4,colframe=teal!35!black,title={Before you move on}]
What does \te{ద్వేషించు} mean? (\textquotedblleft{}To hate\textquotedblright{}.) Say it once more.
\end{tcolorbox}
\section[\texorpdfstring{\te{అసహ్యించుకో} (asahyiñcukō)}{అసహ్యించుకో (asahyiñcukō)}]{\te{అసహ్యించుకో} (asahyiñcukō) — to loathe; to be disgusted by}

\label{lesson:TE-C499-asahyincuko}



\begin{quote}
Before the new one: say the Telugu for to be jealous, then the Telugu for to hate.
\end{quote}

\subsection*{You'll want to know: \te{అసహ్యించుకో}}
\textbf{\te{అసహ్యించుకో}} — \emph{asahyiñcukō} — \textquotedblleft{}to loathe; to be disgusted by\textquotedblright{}.

\begin{quote}
\te{మురికిగా} \te{ఉన్న} \te{వంటగదిని} \te{ఆమె} \te{అసహ్యించుకుంది}.
\end{quote}

\begin{quote}
\emph{murikigā unna vaṇṭagadini āme asahyiñcukundi.}
\end{quote}

\begin{quote}
She was disgusted by the dirty kitchen.
\end{quote}

Past (I …): \te{అసహ్యించుకున్నాను}, \emph{asahyiñcukunnānu}.

The noun is \te{అసహ్యం}, disgust: \te{నాకు} \te{అసహ్యంగా} \te{ఉంది}, it disgusts me.

\begin{cousinweb}
\textbf{\te{అసహ్యించుకో}} (\emph{asahyiñcukō}) is built on Sanskrit \emph{asahya}, \textquotedblleft{}unbearable, not to be endured\textquotedblright{} (Monier-Williams), from \emph{a-}, \textquotedblleft{}not\textquotedblright{}, and \emph{sah}, \textquotedblleft{}to bear\textquotedblright{}; the Telugu endings -\te{ఇంచు} and -\te{కో} make it a verb about one's own feeling.
\end{cousinweb}

\begin{grammarlens}[title={What you've built}]
One more word for describing what you lived and felt.
\end{grammarlens}

\subsection*{Your turn}
\begin{itemize}
\raggedright
  \item \emph{Say it:} \emph{asahyiñcukō}
  \item \emph{Say it:} \emph{asahyiñcukō}, once more
  \item \emph{Say it:} \emph{murikigā unna vaṇṭagadini āme asahyiñcukundi.} — She was disgusted by the dirty kitchen.
  \item \emph{Recall:} say the Telugu for to be jealous, then the Telugu for to hate, then say \emph{asahyiñcukō} again
\end{itemize}

\begin{tcolorbox}[breakable,colback=teal!4,colframe=teal!35!black,title={Before you move on}]
What does \te{అసహ్యించుకో} mean? (\textquotedblleft{}To loathe; to be disgusted by\textquotedblright{}.) Say it once more.
\end{tcolorbox}
\section[\texorpdfstring{\te{పగ} (paga)}{పగ (paga)}]{\te{పగ} (paga) — a grudge; revenge}

\label{lesson:TE-C499-paga}



\begin{quote}
Before the new one: say the Telugu for to hate, then the Telugu for to loathe.
\end{quote}

\subsection*{You'll want to know: \te{పగ}}
\textbf{\te{పగ}} — \emph{paga} — \textquotedblleft{}a grudge; revenge\textquotedblright{}.

\begin{quote}
\te{ఆ} \te{రెండు} \te{కుటుంబాల} \te{మధ్య} \te{పాత} \te{పగ} \te{ఉంది}.
\end{quote}

\begin{quote}
\emph{ā reṇḍu kuṭumbāla madhya pāta paga undi.}
\end{quote}

\begin{quote}
There is an old grudge between those two families.
\end{quote}

\te{పగ} \te{తీర్చుకో}, to take revenge.

\begin{cousinweb}
\textbf{\te{పగ}} (\emph{paga}) is native Dravidian: DEDR 3808 lists Telugu \emph{paga}, \textquotedblleft{}enmity, hatred\textquotedblright{}, and \emph{pagavāḍu}, \textquotedblleft{}an enemy\textquotedblright{}, in the article of Tamil \emph{paku}, \textquotedblleft{}to be split, to be at variance\textquotedblright{}.
\end{cousinweb}

\begin{grammarlens}[title={What you've built}]
One more word for describing what you lived and felt.
\end{grammarlens}

\subsection*{Your turn}
\begin{itemize}
\raggedright
  \item \emph{Say it:} \emph{paga}
  \item \emph{Say it:} \emph{paga}, once more
  \item \emph{Say it:} \emph{ā reṇḍu kuṭumbāla madhya pāta paga undi.} — There is an old grudge between those two families.
  \item \emph{Recall:} say the Telugu for to hate, then the Telugu for to loathe, then say \emph{paga} again
\end{itemize}

\begin{tcolorbox}[breakable,colback=teal!4,colframe=teal!35!black,title={Before you move on}]
What does \te{పగ} mean? (\textquotedblleft{}A grudge; revenge\textquotedblright{}.) Say it once more.
\end{tcolorbox}
\section[\texorpdfstring{\te{చిన్నచూపు} (cinnacūpu)}{చిన్నచూపు (cinnacūpu)}]{\te{చిన్నచూపు} (cinnacūpu) — contempt; looking down on}

\label{lesson:TE-C499-cinnacupu}



\begin{quote}
Before the new one: say the Telugu for to loathe, then the Telugu for a grudge.
\end{quote}

\subsection*{You'll want to know: \te{చిన్నచూపు}}
\textbf{\te{చిన్నచూపు}} — \emph{cinnacūpu} — \textquotedblleft{}contempt; looking down on\textquotedblright{}.

\begin{quote}
\te{పేదవాళ్ళను} \te{చిన్నచూపు} \te{చూడకూడదు}.
\end{quote}

\begin{quote}
\emph{pēdavāḷḷanu cinnacūpu cūḍakūḍadu.}
\end{quote}

\begin{quote}
One must not look down on the poor.
\end{quote}

\te{చిన్నచూపు} \te{చూడు}, to look down on someone: literally a small look.

\begin{cousinweb}
\textbf{\te{చిన్నచూపు}} (\emph{cinnacūpu}) is native Dravidian: \te{చిన్న}, \textquotedblleft{}small, petty\textquotedblright{} (DEDR 2594), and \emph{cūpu}, \textquotedblleft{}sight, a look\textquotedblright{}, the noun of \emph{cūcu}, \textquotedblleft{}to see\textquotedblright{} (DEDR 2735): a look that makes someone small.
\end{cousinweb}

\begin{grammarlens}[title={What you've built}]
One more word for describing what you lived and felt.
\end{grammarlens}

\subsection*{Your turn}
\begin{itemize}
\raggedright
  \item \emph{Say it:} \emph{cinnacūpu}
  \item \emph{Say it:} \emph{cinnacūpu}, once more
  \item \emph{Say it:} \emph{pēdavāḷḷanu cinnacūpu cūḍakūḍadu.} — One must not look down on the poor.
  \item \emph{Recall:} say the Telugu for to loathe, then the Telugu for a grudge, then say \emph{cinnacūpu} again
\end{itemize}

\begin{tcolorbox}[breakable,colback=teal!4,colframe=teal!35!black,title={Before you move on}]
What does \te{చిన్నచూపు} mean? (\textquotedblleft{}Contempt; looking down on\textquotedblright{}.) Say it once more.
\end{tcolorbox}
